% =====================================================================
% PRISMA-ScR Checklist — Supplementary Appendix
% Completed for the scoping review described in the Introduction.
% Based on Tricco et al. (2018), Annals of Internal Medicine.
% 20 essential items + 2 optional items (22 total).
% =====================================================================

\phantomsection
\section*{Multimedia Appendix --- PRISMA-ScR Checklist}\label{sec:appendix-prisma}

This appendix reports the completed Preferred Reporting Items for Systematic reviews and Meta-Analyses extension for Scoping Reviews (PRISMA-ScR) checklist \parencite{Tricco2018_prisma_scr} for the scoping review of multimodal classroom sensing literature described in \Cref{sec:introduction}. The checklist follows the 22-item structure (20 essential, 2 optional) of the PRISMA-ScR statement.

\vspace{1em}

\begin{tabulary}{\textwidth}{@{}p{0.08\textwidth} p{0.42\textwidth} p{0.42\textwidth}@{}}
\toprule
\textbf{Item} & \textbf{PRISMA-ScR Checklist Item} & \textbf{Reported on Page / Section} \\
\midrule
\multicolumn{3}{@{}l}{\textbf{Title}} \\
\cmidrule(l){1-3}
1 & Title: identify the report as a scoping review & The scoping review is identified in \Cref{sec:introduction}, paragraph~2. The present article is a study protocol, not the review itself. \\
\midrule
\multicolumn{3}{@{}l}{\textbf{Abstract}} \\
\cmidrule(l){1-3}
2 & Structured summary: provide a structured summary of the scoping review & Not applicable: the structured abstract reports the protocol study, not the scoping review. \\
\midrule
\multicolumn{3}{@{}l}{\textbf{Introduction}} \\
\cmidrule(l){1-3}
3 & Rationale: describe the rationale for the review in the context of what is already known & \Cref{sec:introduction}, paragraphs~1--2. \\
4 & Objectives: provide an explicit statement of the questions and objectives being addressed with reference to their key elements & \Cref{sec:introduction}, paragraph~2 (identification of the four conjunctive conditions as the review's objective). \\
\midrule
\multicolumn{3}{@{}l}{\textbf{Methods}} \\
\cmidrule(l){1-3}
5 & Protocol and registration: indicate whether a review protocol exists, if and where it can be accessed, and, if available, provide registration information & Not applicable: the scoping review was an internal team synthesis supporting protocol design; no external registration was sought. \\
6 & Eligibility criteria: specify characteristics of the sources of evidence used as eligibility criteria and provide a rationale & \Cref{sec:introduction}, paragraph~2 (three criteria: passive sensing modality, cognitive/attentional outcome, adolescent/young-adult sample in formal education). \\
7 & Information sources: describe all information sources in the search, as well as the date the most recent search was executed & \Cref{sec:introduction}, paragraph~2: Web of Science, Scopus, PubMed, IEEE Xplore, ACM Digital Library; period 2015--2026. \\
8 & Search: present the full electronic search strategy for at least one database, including any limits used & \Cref{sec:introduction}, paragraph~2: three-block search strategy with controlled vocabulary and free-text terms; modality terms, cognitive/attentional construct terms, educational context terms. \\
9 & Selection of sources of evidence: state the process for selecting sources of evidence & \Cref{sec:introduction}, paragraph~2: screening of 1,847 records after deduplication; 287 full texts assessed against eligibility criteria; PRISMA-ScR methodology cited. \\
10 & Data charting process: describe the methods of charting data from the included sources of evidence & \Cref{sec:introduction}, paragraph~2: qualitative synthesis of 22 included studies. \\
11 & Data items: list and define all variables for which data were sought and any assumptions and simplifications made & \Cref{sec:introduction}, paragraph~2: sensing modalities deployed, cognitive/attentional outcome, sample characteristics, and the four conjunctive conditions. \\
12\tablefootnote{Optional item; not performed in this review.} & Critical appraisal of individual sources of evidence: if done, describe the methods used for appraising individual sources of evidence & Not performed: consistent with the scoping-review methodology, which does not require formal risk-of-bias assessment \parencite{Tricco2018_prisma_scr}. \\
13 & Synthesis of results: describe the methods of handling and summarising the data that were charted & \Cref{sec:introduction}, paragraph~2: data synthesised qualitatively by mapping each deployment against the four conjunctive conditions. \\
\midrule
\multicolumn{3}{@{}l}{\textbf{Results}} \\
\cmidrule(l){1-3}
14 & Selection of sources of evidence: give numbers of sources of evidence screened, assessed for eligibility, and included in the review, with reasons for exclusions at each stage & \Cref{sec:introduction}, paragraph~2: 1,847 records screened after deduplication; 287 full texts assessed; 22 studies included in qualitative synthesis. \\
15 & Characteristics of sources of evidence: present characteristics of the included sources of evidence & \Cref{sec:introduction}, paragraph~2: 22 empirical studies (2015--2026). Full characteristics reported in the companion scoping-review article (under review). \\
16\tablefootnote{Optional item; not performed in this review.} & Critical appraisal within sources of evidence: if done, present data on critical appraisal of included sources of evidence & Not performed: formal quality appraisal was not conducted, consistent with the PRISMA-ScR guidance for scoping reviews. \\
17 & Results of individual sources of evidence: for each included source of evidence, present the relevant data that were charted & \Cref{sec:introduction}, paragraph~2: each deployment was classified against the four conjunctive conditions; no single deployment satisfied all four. \\
18 & Synthesis of results: summarise and/or present the charting results as they relate to the review questions and objectives & \Cref{sec:introduction}, paragraphs~2--3: the defining structural limitation (conjunctive gap) is stated, and the present protocol is positioned as addressing it. \\
\midrule
\multicolumn{3}{@{}l}{\textbf{Discussion}} \\
\cmidrule(l){1-3}
19 & Summary of evidence: summarise the main results, including an overview of concepts, themes, and types of evidence available & \Cref{sec:introduction}, paragraphs~2--3. The gap is further discussed in \Cref{subsec:comparison} and \Cref{subsec:principal-findings}. \\
20 & Limitations: discuss the limitations of the scoping review process & \Cref{subsec:limitations}: acknowledged as single-school, single-subject deployment; limitations of the scoping review itself (e.g., exclusion of grey literature) are noted here. Literature published in languages other than English or Spanish was not systematically searched; conference proceedings were hand-searched only for the ACM Digital Library and IEEE Xplore. \\
21 & Conclusions: provide a general interpretation of the results with respect to the review questions and objectives, as well as potential implications and/or next steps & \Cref{sec:introduction}, paragraph~3: the present protocol is designed to fill the conjunctive gap identified by the review. \Cref{subsec:conclusions} reiterates this. \\
\midrule
\multicolumn{3}{@{}l}{\textbf{Funding}} \\
\cmidrule(l){1-3}
22 & Funding: describe sources of funding for the scoping review and the role of funders & The scoping review received no dedicated funding. The broader project is linked to the UCLM teaching-innovation project that developed the ambient platform and the parent national project PID2022-137397NB-I00 (separate scope), as declared in the Acknowledgements. \\
\bottomrule
\end{tabulary}

\endinput
